\documentclass[12pt]{article}
\usepackage[utf8]{inputenc}
\usepackage{bbm}
\usepackage{geometry}
\geometry{a4paper,scale=0.75}
\linespread{1.5}
\usepackage{graphicx} 
\usepackage{float} 
\usepackage{subcaption} 
\usepackage{enumerate}
\usepackage{amsmath}
\usepackage{array}
\newcolumntype{P}[1]{>{\raggedright\arraybackslash}p{#1}}
\usepackage{booktabs}
\usepackage{multirow}
\usepackage{amsfonts}
\usepackage[english]{babel}
\usepackage{amsthm}
\usepackage{dcolumn}
\usepackage{multicol}
\usepackage{stfloats}
\usepackage{placeins}
\usepackage{lscape}
\usepackage[figuresright]{rotating}
\RequirePackage{pdflscape}
\usepackage[toc,page]{appendix}
\usepackage{geometry}
\usepackage{longtable}
\usepackage{comment}
\usepackage{xcolor}
% \usepackage{amsmath}
\usepackage{amssymb}
\usepackage{empheq}
\usepackage{newtxtext,newtxmath}

% -------- enumerated sub-labels (a), (b), … --
\usepackage{enumitem}
\setlist[enumerate,1]{label=(\alph*),ref=\alph*}
% ---------------------------------------------

\usepackage{hyperref}
\hypersetup{hidelinks,
	colorlinks=true,
	allcolors=black,
	pdfstartview=Fit,
	breaklinks=true}
\usepackage{csquotes}
\usepackage{natbib}
\newtheorem{definition}{Definition}
\newtheorem{theorem}{Theorem}
%\newtheorem{proposition}[theorem]{Proposition}
\newtheorem{proposition}{Proposition}[section]
%\newtheorem{lemma}[theorem]{Lemma}
\newtheorem{lemma}{Lemma}[section]
%\newtheorem{corollary}[theorem]{Corollary}
\newtheorem{corollary}{Corollary}[section]
\newtheorem*{remark}{Remark}
\newtheorem{example}{Example}
\newtheorem{exercise}{Exercise}[section]
\numberwithin{equation}{section}
\newtheorem{assumption}{Assumption}[section] % number within sections

\title{Notes on \cite{bloomReallyUncertainBusiness2018}}
\author{Harlly Zhou \\ \texttt{hzhouci@ust.hk}}
\date{\today}

\begin{document}
\maketitle

\paragraph{Motivation} Uncertainty is importqnt. Try to evaluate the role of uncertainty for business cycle.

\paragraph{Measurement} A firm produces output according to
\begin{align*}
    y_{j,t} = A_t z_{j,t} f(k_{j,t}, n_{j,t})
\end{align*}
where $A_t$ is the agggregate productivity and $z_{j,t}$ is the idiosyncratic component. They follow an $AR(1)$ process:
\begin{align*}
    \log A_t &= \rho^A \log A_{t-1} + \sigma_{t-1}^A \epsilon_t\\
    \log z_{j,t} &= \rho^z \log z_{j,t-1} + \sigma_{t-1}^z \epsilon_{j,t}.
\end{align*}
Let the variances move over time according to two-state Markov chains.

\paragraph{Empirics} Uncertainty is countercyclical, but the causality seems to be two-way.
\begin{enumerate}
    \item Dispersion in plant-level innovation is countercyclical. 
    \begin{enumerate}
        \item Data: Census panel of manufacturing establishments. 
        \item Calculate the establishment-lebel TFP $\hat{z}_{j,t}$ according to Foster, Haltiwanger and Krizan (2000). Define TFP \textbf{shock} as $e_{j,t}$:
        \begin{align*}
            \log \hat{z}_{j,t} = \rho \log \hat{z}_{j,t-1} + \underbrace{\mu_j}_{\text{establishment FE}} + \underbrace{\lambda_t}_{\text{time FE}} + e_{j,t}.
        \end{align*}
        Define microeconomic uncertainty $\sigma^{\hat z}_{t-1}$ as the cross-sectional dispersion of $e_{j,t}$ calculated on an annual basis.
        \item Empirical strategy: Regress dispersion measurements (standard deviation, skewness, kurtosis, and IQR\footnote{Interquartile range, which is the upper quartiel minus the lower quartile.}) on the number of recession quarters in a year (\textit{e.g.}, recession only in Q4 will be assigned with number 0.25).
    \end{enumerate}
    \item Industry-level uncertainty is countercyclical.
    \item Macroeconomic measures of uncertainty: Estimated using GARCH(1,1) model on Basu, Fernald and Kimball (2006) (seems to have been mentioned in Byoungchan's class).
\end{enumerate}


\bibliographystyle{../draft/aer}
\bibliography{../draft/network_uncertainty}

\end{document}
